	\paragraph{Kết luận:}
	Dưới tác động của hiện tượng phương sai thay đổi, sai số chuẩn tính toán từ thuật toán OLS thông thường có xu hướng bị hạ thấp quá mức so với thực tế (lạc quan sai lệch). Sau khi hiệu chỉnh bằng chuẩn vững HC3, các sai số chuẩn (\textit{Std. Error}) đều tăng lên. Tuy nhiên, mức độ ý nghĩa thống kê của hai biến độc lập cốt lõi là \texttt{income} và \texttt{education} \textbf{vẫn giữ nguyên vẹn ký hiệu ý nghĩa tối cao ($***$)}. Điều này chứng minh mối quan hệ tác động tích cực của thu nhập và trình độ học vấn đến uy tín nghề nghiệp là hoàn toàn bền vững và có độ tin cậy khoa học rất cao.
	\begin{itemize}
		\item \textbf{B-P Test:} Nhanh, mạnh khi phương sai thay đổi tuyến tính, nhưng dễ sai lệch nếu phân phối sai số bất thường (Khắc phục bằng Koenker studentized).
		\item \textbf{White Test:} Toàn diện, không cần định hình trước cấu trúc, nhưng hao tốn quá nhiều bậc tự do, yếu ở mẫu nhỏ.
		\item \textbf{Robust SE (HC3):} Không phải là phép kiểm định mà là "liều thuốc chữa bệnh". Giữ vững tính không chệch của OLS, hiệu chỉnh lại $SE$ để đưa ra quyết định thống kê khoa học nhất.
	\end{itemize}
